For fold \(e\), condition on the opposite-fold sigma-field \(\mathscr F_{-e}\). By \(\mathrm{ass{:}two\text{-}arm\text{-}sampling}\), the evaluation observations are independent of \(\mathscr F_{-e}\). With the notation of the pilot lemma, the exact product expansion is
\[
\widehat U_e=U_e^{\mathrm{can}}(\widehat B)
+2\bar Z_{1,e}^\top\widehat Bd-2\bar Z_{0,e}^\top\widehat Bd+d^\top\widehat Bd. \tag{1}
\]
Apply the inverse-cdf Taylor identity (12) of that lemma to the conditional transport functional, and use linearity in the common marginal. The plug-in is the constant term, the displayed scores together with the two linear terms in (1) are the first Hoeffding projection, and \(d^\top\widehat Bd\) is the conditional mean of the quadratic correction. After averaging the two folds and separating the projection complement, this gives the exact decomposition
\[
\widehat\theta^{\mathrm{an}}_{J_y,J_z,n}-\theta_n
=\mathbb L_{J_y,J_z,n}+\mathbb U_{J_y,J_z,n}
+B^y_{J_y,n}+B^z_{J_z,n}+R_{J_y,J_z,n}.                              \tag{2}
\]
Here \(\mathbb U\) is the average of \(U_e^{\mathrm{can}}(B)\), while \(\mathbb L\) is the oracle first projection. Equations (13)--(20) of \(\mathrm{lem{:}observable\text{-}pilot\text{-}control}\) give, uniformly,
\[
R_{J_y,J_z,n}
=O_P\!\left\{\varepsilon^{2+a_y}+\rho_n^{2a_y/(2+3/s_y+d/s_z)}\varepsilon^2\right\}
+o_P\!\left\{\frac{\rho_n}{\sqrt n}+\frac{A_{J_y,J_z,n}}n\right\}. \tag{3}
\]
In particular, (3) includes the conditional linear standard deviation explicitly computed there; it is absorbed because \(\sqrt n\varepsilon\to\infty\).

Let \(H_z(y)=\int_0^y\Delta_n(v,z)dv\) and \(u_z=\mathcal G_{n,z}\Delta_n(\cdot,z)\). Integrating the Neumann equation and using zero conditional mass gives
\[
-\bar q_n(y,z)u_z'(y)=H_z(y),
\qquad
\int_0^1\Delta_nu_z\,dy=\int_0^1\frac{H_z(y)^2}{\bar q_n(y,z)}\,dy.       \tag{4}
\]
Implicit differentiation of the two conditional quantile maps shows that the squared \(L^2(P_{0,n})\oplus L^2(P_{1,n})\)-norm of the centered oracle score is bounded above and below by fixed multiples of the right side of (4), integrated against \(R_n\). The density bounds and the definition of \(\rho_n\) therefore imply
\[
V_{L,n}=n\operatorname{Var}(\mathbb L_{J_y,J_z,n})\asymp\rho_n^2,
\qquad \operatorname{sd}(\mathbb L_{J_y,J_z,n})\asymp\rho_n/\sqrt n. \tag{5}
\]
At equality \(H_z=0\); both the value derivative and \(\mathbb L\) vanish.

For the canonical projection, independence and diagonal deletion leave only equal or reversed index pairs. With \(m=n/2\), direct counting gives
\[
\begin{split}
\operatorname{Var}(\mathbb U)
={}&\frac{\operatorname{tr}\{(B\Gamma_0)^2\}+\operatorname{tr}\{(B\Gamma_1)^2\}}{m(m-1)}
+\frac{2\operatorname{tr}(B\Gamma_0B\Gamma_1)}{m^2}\\
={}&\frac{4\operatorname{tr}\{(\Gamma^{1/2}B\Gamma^{1/2})^2\}}{n^2}\{1+o(1)\}
=\frac{A_{J_y,J_z,n}^2}{n^2}\{1+o(1)\}.                            \tag{6}
\end{split}
\]
The harmless \(1+o(1)\) is absorbed by the normalization convention in the statement; equivalently the displayed standard-deviation equality holds for its canonical finite-sample normalization, and the raw pair count differs by \(O(m^{-1})\).

The corrected spectrum proposition supplies
\[
|B^y_{J_y,n}|\lesssim2^{-2(s_y+1)J_y},
\qquad |B^z_{J_z,n}|\lesssim2^{-2s_zJ_z}.                              \tag{7}
\]
The theorem's premise turns (3) into little-oh of the sum of the four fluctuation and bias scales. Substitution into (2), together with (5)--(7), proves the expansion.

For the root-mean-square assertion, tail integration in (2) of the pilot proof gives fixed fourth moments for the pilot errors. Applying the deterministic bounds (13)--(15), the conditional variance identities (16)--(17), and then integrating over the pilots yields
\[
\|R_{J_y,J_z,n}\|_2
\le C\{\varepsilon^{2+a_y}+\rho_n^{2a_y/(2+3/s_y+d/s_z)}\varepsilon^2\}
+o\!\left\{\frac{\rho_n}{\sqrt n}+\frac A n\right\}.             \tag{8}
\]
This is the required uniform \(L^2\) upgrade. At the declared tuning,
\[
\frac{A_{J_y^*,J_z^*,n}}n
\asymp\frac{2^{dJ_z^*/2}}n
\asymp2^{-2s_zJ_z^*}
\asymp2^{-2(s_y+1)J_y^*}
\asymp r_{n,d}(s_z).                                                    \tag{9}
\]
Minkowski's inequality applied to (2), using (5)--(9), proves the displayed root-mean-square scale and its absorption sub-regime. No empirical covariance occurs in (1)--(9), so covariance estimation is not needed for the point estimator.